\section{Real Data Application: Unspanned Macroeconomic Signals in Bond Risk Premia}
\label{sec:app}

We ask whether a large macroeconomic panel contains stable predictive
information for U.S.\ Treasury excess returns beyond that already summarized by
the yield curve. For a compact main presentation, we report two- and three-year
bonds here and the corresponding four- and five-year results in
Appendix~\ref{app:long_maturity_results}; all maturities use the same forecast
sample, information set, and specification.

\subsection{Data, forecasting design, and methods}
\label{sec:app_design_v7}

The targets are one-year log excess holding returns on two- through five-year
Fama--Bliss zero-coupon bonds \citep{LudvigsonNg2009}. The January 1970--December
2023 sample contains 118 complete series from the November 2024 FRED-MD vintage
\citep{McCrackenNg2016}, transformed by the database codes and standardized
within each estimation window.

At each forecast origin $\tau$, a 400-month rolling window supplies 388 fully
realized pairs $(X_s,r^{(n)}_{s\to s+12})$ with $s+12\leq\tau$. Every estimation
step is repeated at the origin, producing 236 common forecasts from May 2004
through December 2023.

Every specification includes three yield-curve principal components (level,
slope, and curvature) \citep{LittermanScheinkman1991} and the
\citet{CochranePiazzesi2005} (CP) return-forecasting factor. Their loadings and
coefficients are re-estimated using only information available at each origin,
so selected macro variables must add to a level--slope--curvature--CP benchmark.

We compare recycled FAR-Knockoff with \textsc{Lasso}, \textsc{FarmSelect}
\citep{FanKeWang2020}, \textsc{Model-X}, and \textsc{Fixed-X} knockoffs
\citep{BarberCandes2015,CandesFanEtal2018}, using identical observations,
controls, and post-selection OLS refits. FAR-Knockoff and FarmSelect choose one
to six macro factors by their implemented information criterion. For prediction,
FAR-Knockoff uses $q=0.2$ and averages its LCD statistics over 20 random
original--knockoff relabelings to reduce sensitivity to a single labeling.
Appendix~\ref{app:long_maturity_results} reports factor-cap sensitivity.

\subsection{Short-end forecast performance}
\label{sec:app_results_v7}

Table~\ref{tab:empirical_comparison_v7} reports RMSE; parentheses give mean
selected macro-model size among nonempty windows, excluding the four controls.
The final rows test FAR-Knockoff against the nested benchmark.

\begin{table}[H]
\caption{Out-of-sample forecast performance at the short end}
\label{tab:empirical_comparison_v7}
\centering
\small
\setlength{\tabcolsep}{8pt}
\renewcommand{\arraystretch}{1.08}
\begin{tabular}{lcc}
\toprule
 & 2 years & 3 years \\
\midrule
Yield factors + CP       & 1.458 (--) & 2.812 (--) \\
\textsc{Lasso}           & 2.516 (23.33) & 4.638 (18.81) \\
\textsc{FarmSelect}      & 2.988 (43.81) & 5.322 (41.49) \\
\textsc{Model-X}         & 1.911 (3.69) & 3.042 (3.28) \\
\textsc{Fixed-X}         & 1.920 (1.74) & 3.208 (2.62) \\
\textsc{FAR-Knockoff}    & \textbf{1.380} (2.78) & \textbf{2.763} (2.10) \\
\midrule
Clark--West statistic    & 2.49 & 1.78 \\
One-sided $p$-value      & 0.0065 & 0.0377 \\
\bottomrule
\end{tabular}

\vspace{2pt}
\begin{minipage}{0.92\textwidth}
\footnotesize
\emph{Notes:} The first six rows report RMSE, with conditional mean selected
size in parentheses. RMSE and tests use all 236 windows, including those with
no selected macro variables. Clark--West tests compare FAR-Knockoff with yield
factors + CP using Bartlett HAC standard errors with lag 11.
\end{minipage}
\end{table}

FAR-Knockoff has the lowest RMSE at both short maturities while retaining only
two to three macro variables in nonempty windows. At two years, RMSE falls
from 1.458 to 1.380, a 10.5\% reduction in mean squared forecast error;
the one-sided Clark--West $p$-value is 0.0065. At three years, the corresponding
reduction is 3.5\% ($p=0.0377$), but it is sensitive to the macro factor cap.

This maturity localization has a close precedent in the macro-finance
literature. \citet{LudvigsonNg2009} document macro-factor predictability across
two- to five-year bonds, whereas the robust reanalysis of
\citet{BauerHamilton2018} finds substantially weaker evidence beyond the yield
curve and identifies the two-year bond as the exceptional maturity in the
relevant spanning exercises. Our rolling forecasts are consistent with a
maturity-specific increment beyond a strong yield benchmark. The selected
variables below suggest a connection to policy and cyclical news, although
forecasting associations alone do not identify a causal channel.

The cross-method comparison shows that the gain does not arise from adding many
predictors. \textsc{Lasso} and \textsc{FarmSelect} select much larger models
and substantially increase forecast loss. \textsc{Model-X} and
\textsc{Fixed-X} are more parsimonious but do not improve on the benchmark.
FAR-Knockoff combines the smallest short-end forecast loss with a sparse
selected block.

\subsection{Selected predictors and economic interpretation}
\label{sec:app_interpretation_v7}

FAR-Knockoff selects coordinates of the estimated idiosyncratic component
$\widehat U$, after common macro factors and the always-included yield controls
have been partialled out. The equivalent representation in Section~2 retains
each original predictor's coefficient, so a selected series identifies its
incremental information beyond the factors and controls, not a structural or
causal shock. Table~\ref{tab:empirical_selection_v7} reports frequencies over
all 236 windows.

\begin{table}[H]
\caption{Most frequently selected short-end residual macro predictors}
\label{tab:empirical_selection_v7}
\centering
\footnotesize
\setlength{\tabcolsep}{4.8pt}
\renewcommand{\arraystretch}{1.0}
\begin{tabular}{clrl}
\toprule
Maturity & FRED-MD series & Frequency & Incremental economic channel \\
\midrule
2 years & \texttt{T5YFFM} & 24.6\% & policy-path / term-premium news \\
        & \texttt{PERMITNE} & 24.6\% & forward-looking housing activity \\
        & \texttt{TB3SMFFM} & 13.1\% & short-rate spread / policy expectations \\
3 years & \texttt{T5YFFM} & 26.3\% & policy-path / term-premium news \\
        & \texttt{VIXCLSx} & 17.4\% & uncertainty and risk-bearing capacity \\
        & \texttt{CES0600000007} & 16.1\% & goods-producing workweek \\
\bottomrule
\end{tabular}
\end{table}

\enlargethispage{\baselineskip}
The two-year selections suggest a policy-expectations and interest-sensitive
activity channel. \texttt{T5YFFM} and \texttt{TB3SMFFM} are the five-year
Treasury and three-month Treasury-bill spreads over the federal funds rate.
Their residual components describe relative rate movements not absorbed by
the macro factors and yield controls. For a one-year holding return on a
two-year bond, revisions to the short end of the yield curve are a natural
source of price variation. \texttt{PERMITNE}, Northeast housing permits,
adds forward-looking information about rate-sensitive demand. Its residual
is unusual housing activity relative to the common macro cycle, rather than
simply another measure of that cycle. The recurring combination of rate
spreads and housing activity is therefore consistent with policy-transmission
news beyond the yield benchmark, without establishing a structural mechanism.

At three years, \texttt{T5YFFM} remains prominent, while \texttt{VIXCLSx}
and \texttt{CES0600000007}, the goods-producing workweek, indicate uncertainty
and labor-utilization channels. The latter is consistent with the real-activity
component of bond predictability in \citet{LudvigsonNg2009} and models in which
output risks affect bond risk premia without being fully spanned by yields
\citep{JoslinPriebschSingleton2014}. The VIX association is also consistent
with evidence linking bond-return predictability to macroeconomic conditions
and short-rate expectations \citep{GarganoPettenuzzoTimmermann2019}.
These interpretations concern residual predictive information, not the signs
of causal effects, which selection frequencies do not identify.
